% ============================================================================
% Fichier  : partie5/P05_C17_continuite.tex
% Titre    : Continuité d'Activité : PCA et PRA
% Auteur   : MINKA MI NGUIDJOÏ Thierry Emmanuel
% ============================================================================
% Sources : ISO 22301:2019 ; ISO/IEC 27002:2022. Mairie d'Eséka : cas fictif.
% ============================================================================

\chapter{Continuité d'Activité : PCA et PRA}
\label{chap:continuite}
\chaptermark{PCA et PRA}

\epigraph{%
    « Une sauvegarde que l'on n'a jamais restaurée\\
    n'est pas une sauvegarde~: c'est un espoir. »%
}{Principe du cours}

\niveauChapitre{Intermédiaire}{Le Chapitre~\ref{chap:technique}.}

\begin{boxobjectifs}
\begin{itemize}
    \item Distinguer plan de continuité (\sigle{PCA}) et plan de reprise
          (\sigle{PRA}).
    \item Conduire une analyse d'impact sur l'activité (\sigle{BIA}).
    \item Définir \sigle{RTO}, \sigle{RPO}, \sigle{MTTR}, \sigle{MTBF}.
    \item Appliquer la règle de sauvegarde \mbox{3-2-1} et tester les
          restaurations.
    \item Construire une stratégie de continuité adaptée à la mairie.
\end{itemize}
\end{boxobjectifs}

\bigskip

% ============================================================================
\section*{Question directrice : comment tenir le service quand le pire arrive ?}
% ============================================================================

L'incident a montré qu'une panne peut tout arrêter. La continuité prépare la
mairie à tenir et à reprendre. Ce chapitre construit le \sigle{PCA} et le
\sigle{PRA}.

\medskip
\textbf{Plan du chapitre.}
\begin{enumerate}
    \item PCA et PRA (section~\ref{sec:c17-pcapra}).
    \item L'analyse d'impact (section~\ref{sec:c17-bia}).
    \item Les indicateurs de continuité (section~\ref{sec:c17-indic}).
    \item Sauvegarde et restauration (section~\ref{sec:c17-sauvegarde}).
    \item La stratégie de la mairie (section~\ref{sec:c17-strategie}).
\end{enumerate}

\separateur

% ============================================================================
\section{PCA et PRA}
\label{sec:c17-pcapra}
% ============================================================================

\begin{boxdef}
\termecle{Plan de continuité d'activité} (\sigle{PCA})~: dispositif qui maintient
les activités critiques pendant un sinistre, éventuellement en mode dégradé.
\termecle{Plan de reprise d'activité} (\sigle{PRA})~: dispositif qui rétablit le
système après le sinistre. Le \sigle{PCA} tient~; le \sigle{PRA} répare.
\end{boxdef}

Pour la mairie, le \sigle{PCA} permet de continuer à délivrer des actes (même
sur registre papier de secours)~; le \sigle{PRA} remet en service le système
numérique.

% ============================================================================
\section{L'Analyse d'Impact}
\label{sec:c17-bia}
% ============================================================================

\begin{boxdef}
\termecle{Analyse d'impact sur l'activité} (\sigle{BIA})~: étude qui identifie
les activités critiques, évalue l'impact de leur interruption et fixe les délais
tolérables. Elle fonde toute la stratégie de continuité.
\end{boxdef}

Le \sigle{BIA} répond à~: quelles activités ne peuvent s'arrêter~? combien de
temps peut-on tenir sans elles~? quelle perte de données est tolérable~? À
Eséka, l'état civil est l'activité la plus critique.

% ============================================================================
\section{Les Indicateurs de Continuité}
\label{sec:c17-indic}
% ============================================================================

\begin{tableaucours}{L{2.0cm}L{10.2cm}}{%
    \tableauHeader{Indicateur} &
    \tableauHeader{Définition}
}
RTO & Délai maximal de remise en service visé \\
RPO & Perte de données maximale acceptable (point de reprise) \\
MTTR & Temps moyen de réparation \\
MTBF & Temps moyen entre deux pannes \\
\end{tableaucours}
\vspace{-0.8em}
\begin{center}{\small\itshape Tableau~17.1. Indicateurs de continuité}\end{center}

\sigle{RTO} et \sigle{RPO} commandent la stratégie~: un \sigle{RPO} d'une heure
impose des sauvegardes au moins horaires~; un \sigle{RTO} de quatre heures
impose des moyens de reprise rapides.

% ============================================================================
\section{Sauvegarde et Restauration}
\label{sec:c17-sauvegarde}
% ============================================================================

\begin{boxdef}
La règle \termecle{3-2-1}~: conserver \textbf{3} copies des données, sur
\textbf{2} supports différents, dont \textbf{1} hors site. Elle protège contre
la panne, l'erreur et le sinistre local.
\end{boxdef}

\begin{boxprocedure}
\textbf{Constatation.} Avant la démarche, la mairie avait une seule copie de
l'état civil, sur le même serveur que l'original, jamais testée.

\medskip
\textbf{Analyse} (confiance élevée). Cette « sauvegarde » violait la règle
3-2-1 (une copie, un support, sur site) et n'avait jamais été restaurée~: elle
ne protégeait rien. La mesure corrective impose trois copies, dont une hors
site, et un \emph{test de restauration} mensuel. Une sauvegarde non testée est
réputée inexistante.
\end{boxprocedure}

\begin{boxmetier}{DSI}
\textbf{Tester, pas seulement sauvegarder.} Le \sigle{DSI} planifie des tests
de restauration réguliers~: c'est le seul moyen de savoir que la sauvegarde
fonctionne \emph{avant} d'en avoir besoin. Le jour du sinistre est le pire
moment pour découvrir qu'elle est illisible.
\end{boxmetier}

% ============================================================================
\section{La Stratégie de la Mairie}
\label{sec:c17-strategie}
% ============================================================================

La mairie d'Eséka adopte une stratégie proportionnée~: sauvegarde quotidienne
de l'état civil selon la règle 3-2-1, copie hors site chez la tutelle, registre
papier de secours pour le mode dégradé, et test de restauration mensuel. Le
\sigle{PCA} et le \sigle{PRA} sont écrits, diffusés et exercés.

\separateur

\begin{boxresume}
\textbf{Ce qu'il faut retenir de ce chapitre~:}
\begin{itemize}
    \item Le \textbf{PCA} maintient l'activité~; le \textbf{PRA} rétablit le
          système.
    \item Le \textbf{BIA} identifie les activités critiques et fixe les délais
          tolérables.
    \item \textbf{RTO} et \textbf{RPO} commandent la stratégie de continuité.
    \item La règle \textbf{3-2-1} et le \textbf{test de restauration} font
          d'une copie une vraie sauvegarde.
    \item La mairie adopte une stratégie \textbf{proportionnée} et
          \textbf{exercée}.
\end{itemize}
\end{boxresume}

\begin{boxinfo}
\textbf{Points de contrôle de l'élève.} Je sais~: distinguer PCA et PRA~;
conduire un BIA~; définir RTO, RPO, MTTR, MTBF~; appliquer la règle 3-2-1~;
bâtir une stratégie de continuité.
\end{boxinfo}

\bigskip

\begin{boxexo}[title={\ding{45}\quad Exercice~17.1. PCA ou PRA \niveaubase}]
« Continuer à délivrer des actes pendant la panne » relève-t-il du PCA ou du
PRA~?
\end{boxexo}

\begin{boxexo}[title={\ding{45}\quad Exercice~17.2. Les indicateurs \niveaubase}]
Définissez RTO et RPO.
\end{boxexo}

\begin{boxexo}[title={\ding{45}\quad Exercice~17.3. La règle 3-2-1 \niveaubase}]
Énoncez la règle 3-2-1 de sauvegarde.
\end{boxexo}

\begin{boxexo}[title={\ding{45}\quad Exercice~17.4. BIA \niveaumoyen}]
Quelle est l'activité la plus critique de la mairie, et pourquoi~?
\end{boxexo}

\begin{boxexo}[title={\ding{45}\quad Exercice~17.5. RPO et fréquence \niveaumoyen}]
Un RPO d'une heure impose quelle fréquence de sauvegarde~? Justifiez.
\end{boxexo}

\begin{boxexo}[title={\ding{45}\quad Exercice~17.6. Fausse sauvegarde \niveaumoyen}]
Pourquoi la « sauvegarde » initiale de la mairie ne protégeait-elle rien~?
\end{boxexo}

\begin{boxexo}[title={\ding{45}\quad Exercice~17.7. Tester \niveaumoyen}]
Pourquoi faut-il tester les restaurations, et à quelle fréquence~?
\end{boxexo}

\begin{boxexo}[title={\ding{45}\quad Exercice~17.8. Mode dégradé \niveaumoyen}]
Décrivez un mode dégradé pour l'état civil en cas de panne prolongée.
\end{boxexo}

\begin{boxexo}[title={\ding{45}\quad Exercice~17.9. Stratégie 3-2-1 \niveauexpert}]
Concevez la stratégie de sauvegarde 3-2-1 complète de la mairie d'Eséka.
\end{boxexo}

\begin{boxexo}[title={\ding{45}\quad Exercice~17.10. PRA rédigé \niveauexpert}]
Rédigez un plan de reprise d'activité d'une page pour le serveur d'état civil
(RTO, RPO, étapes, responsables).
\end{boxexo}

\bigskip

\begin{boxinfo}
\textbf{Pour aller plus loin.}
\begin{itemize}
    \item ISO 22301~\cite{iso22301}~; ISO/IEC 27002~\cite{iso27002}.
\end{itemize}
\end{boxinfo}

\begin{boxinfo}
\textbf{Transition.} Les mesures en place, il faut les piloter~: la partie VI
s'ouvre au Chapitre~\ref{chap:indicateurs} sur la surveillance et les
indicateurs.
\end{boxinfo}
